\chapter{Agentic Evolutionary Frameworks}
\label{chapter:EvolutionaryFrameworks}
% EVOLVE-BLOCK-START

The pivot from orchestration to evolution is a single conceptual step.\index{Evolutionary framework}
Across the foundations half of the course, a loop was assembled: generate candidates, evaluate them with the shared harness, select the best, and repeat.
The generate–evaluate–select (GES) loop of Chapter~\ref{chapter:Orchestration} is already an optimization procedure, a search over the space of candidate models driven by the evaluation score.

One ingredient distinguishes a random search from an evolutionary search: \emph{mutation}.\index{Mutation operator}
A mutation operator derives new candidates from survivors by making structured, small edits to them.
Survivors that score well carry structure worth preserving; mutation explores variations on that structure.
Over generations, the population improves because good structure is inherited and refined.

This chapter formalizes the evolutionary loop, develops the components that make it work, and traces the lineage from classical genetic programming through FunSearch and AlphaEvolve to the systems a researcher builds in the capstone.

\begin{remark}
It is worth asking why an \emph{external} loop is needed at all, now that a capable model can reason at length before it answers: iterating, reconsidering, and backtracking within a single call (Chapter~\ref{chapter:AgentLoop}).
For a problem whose solution lies inside the space the model can search internally, that deliberation may indeed do the work, and it is worth trying as a baseline before building a harness.
Two things keep the loop irreducible when it is.
First, the \emph{evaluator is an external oracle}: the model can reason \emph{about} whether a candidate works, but the ground-truth signal (does this program run, what error does the simulator return, what does the measurement say) comes from executing against the world, not from further thought (Spine~1 of Section~\ref{section:TwoSpines}, in its starkest form).
Second, search over a large, rugged space of structured artifacts is a different regime from bounded per-call reasoning: the population, its memory of what has already failed, and the selection pressure across generations explore a landscape no single forward pass exhausts.
Even a model that searches \emph{within} its own deliberation --- as in Tree-of-Thoughts (Yao et al., 2023), which explores intermediate reasoning steps for a single problem instance --- is doing internal, self-scored search, not the external, harness-scored search across a population this chapter builds.
Extended thinking makes each \emph{proposal} smarter; it does not replace the external check or the search around it.
\end{remark}

\section{The Evolutionary Loop}
\label{section:EvolutionaryLoop}

Definition~\ref{definition:GenEvalSel} is already in place; this chapter adds the one missing ingredient, a mutation operator, and shows what changes when candidates are derived from survivors rather than drawn afresh.

\begin{definition}
\label{definition:EvolutionaryLoop}
An \emph{agentic evolutionary framework} is a tuple $(\Artifacts, \Fitness, \Mutate, \Select, \Population_0)$ where
\begin{itemize}
\item $\Artifacts$ is the space of artifacts (expressions, programs, designs),
\item $\Fitness : \Artifacts \to \RealNumbers$ is the fitness function (lower is better, following the convention of the course harness),
\item $\Mutate : \Artifacts \to \Artifacts$ is the mutation operator that derives a variant from a survivor,
\item $\Select : 2^{\Artifacts} \to 2^{\Artifacts}$ is the selection rule that retains the best members of a population, and
\item $\Population_0 \subset \Artifacts$ is the initial population (the set of candidates before any evolution begins).
\end{itemize}
Starting from $\Population_0$, the loop produces generations $\Population_1, \Population_2, \ldots$ according to
\begin{equation}
\Population_{t+1} = \Select(\Population_t \cup \{\Mutate(a) : a \in \Select(\Population_t)\}) .
\label{equation:EvolutionaryLoop}
\end{equation}
\end{definition}

In words, the tuple names the five things that must be supplied before any evolution can run, and nothing else.
$\Artifacts$ says what kind of object is being evolved (here, an algebraic expression such as \code{x0 * exp(-x1)}).
$\Fitness$ says how good one such object is, as a single number a machine can compute without asking anyone's opinion.
$\Mutate$ says how to make a new object out of an existing one.
$\Select$ says which objects survive a round.
$\Population_0$ says where the search starts.
A biologist will recognize the vocabulary immediately: $\Artifacts$ is the space of possible genotypes, $\Fitness$ is the environment that scores them, $\Mutate$ is the source of heritable variation, and $\Select$ is the pressure deciding which individuals get to reproduce.
The analogy is close enough to be load-bearing throughout this chapter, with one difference worth naming up front: the variation operator here is not random, it is a language model carrying priors about which variations are worth trying at all (Section~\ref{section:MutationOperators}).

The recurrence uses $\Select$ twice, and the two uses do different jobs.
The inner $\Select(\Population_t)$ chooses the \emph{parents}, the survivors that will be mutated; the outer $\Select$ prunes the enlarged pool back down to a next generation of fixed size.
Reading Eq.~\eqref{equation:EvolutionaryLoop} from the inside out therefore traces one round in four moves: choose parents, mutate each parent once, pool parents together with their offspring, keep the best of that pool.

\begin{example}
The tuple becomes concrete as soon as it is filled in.
For the symbolic-regression runs of this chapter, $\Artifacts$ is the set of algebraic expressions built from the variables \code{x0, x1, x2}, numeric constants, and a fixed operator set ($+$, $-$, $\times$, $\div$, \code{exp}, \code{log}, \code{sqrt}), capped at some node budget so expressions cannot grow without limit.
$\Fitness$ is the composite score $\Score = \mathrm{MSE} + \lambda C$ returned by the evaluation harness of Chapter~\ref{chapter:Evaluation}.
$\Mutate$ is one call to a language model with the prompt of Section~\ref{section:MutationOperators}, returning one edited expression.
$\Select$ keeps the eight lowest-scoring candidates.
$\Population_0$ is the twenty expressions a first proposer pass produced before any evolution began.
Nothing in that list is exotic, and every entry is a knob: changing the operator set, the score, the prompt, the survivor count, or the seeding pass changes the search, which is why the capstone is mostly a sequence of deliberate choices about these five slots.
\end{example}

The key insight from~\eqref{equation:EvolutionaryLoop} is that the next generation is drawn from the union of (a) the survivors of the current generation (the \emph{elite}), and (b) mutations of those survivors.
Good structure is not discarded, it is carried forward and explored around.
Figure~\ref{figure:EvolutionaryLoop} draws the loop, with \emph{Mutate} shown as a node distinct from \emph{Generate}: mutation is the single operator absent from the generate--evaluate--select loop of Chapter~\ref{chapter:Orchestration}, and it is what makes the candidates depend on the survivors.

\begin{figure}[htb!]
\begin{center}
%% evolutionary-loop.tex: the agentic evolutionary loop (Chapter 10).
%% House conceptual style; Mutate is drawn as a node distinct from Generate, since it
%% is the one operator absent from the generate-evaluate-select loop of Chapter 9.
%% _concept-style.tex: shared TikZ styles for the course's CONCEPTUAL block diagrams.
%% The leading underscore marks this as the one file in figures/ that is not a
%% figure; every other figures/*.tex holds exactly one figure body.
%%
%% \input this at the top of every conceptual figure (before \begin{tikzpicture})
%% so each such diagram speaks one visual language: rounded "block" nodes, stealth
%% arrows, and natural `<hue>!15' fills. Redefining these styles on each \input is
%% idempotent and harmless. The reference for the house parameters is
%% resources/STYLE-figures.md. Figure~\ref{figure:AgentLoop} is the exemplar.
%%
%% Scope: use for conceptual/relational diagrams (boxes-and-arrows). Do NOT use for
%% product mock-ups (the rig figures have their own shared language) or for
%% data/plot figures.
\tikzset{
  % the standard block node: geometry only — apply a `fill=<hue>!15' per node.
  conceptbox/.style={draw, rounded corners=4pt, minimum width=2.4cm,
                     minimum height=0.85cm, align=center, font=\small},
  % a flow arrow between blocks.
  conceptflow/.style={->, >=stealth', thick},
  % an edge/annotation label.
  conceptlabel/.style={font=\scriptsize, align=center},
}%

\begin{tikzpicture}
% One-time seed on the left; the steady-state cycle runs clockwise on the right.
\node[conceptbox, fill=yellow!20, minimum width=3.4cm] (pop) at (0, 2.0) {Initial population $\Population_0$};
\node[conceptbox, fill=blue!8]    (gen) at (0,    0)    {Generate};
\node[conceptbox, fill=green!20]  (fit) at (5.2,  0)    {Evaluate\\(Fitness)};
\node[conceptbox, fill=orange!20] (sel) at (10.4, 0)    {Select};
\node[conceptbox, fill=blue!15]   (mut) at (5.2, -2.2)  {Mutate};

\draw[conceptflow] (pop) -- node[conceptlabel, midway, right]{seed}    (gen);
\draw[conceptflow] (gen) -- node[conceptlabel, midway, above]{initial} (fit);
\draw[conceptflow] (fit) -- node[conceptlabel, midway, above]{scored}  (sel);
\draw[conceptflow] (sel.south) -- (sel.south |- mut.east)
    -- node[conceptlabel, midway, below]{survivors $\Population_{t+1}$} (mut.east);
\draw[conceptflow] (mut) -- node[conceptlabel, midway, right]{candidates} (fit);
\end{tikzpicture}%

\caption{The agentic evolutionary loop.
An initial population seeds the \emph{Generate} step once; thereafter the steady-state cycle is evaluate~$\to$~select~$\to$~\emph{mutate}, and it is \emph{Mutate}, not \emph{Generate}, that takes the survivors as input, so each round's candidates are derived from the previous round's elite rather than drawn afresh.
Mutation is the one node absent from the generate--evaluate--select loop of Chapter~\ref{chapter:Orchestration}; the fitness function is the evaluation harness of Chapter~\ref{chapter:Evaluation}.}
\label{figure:EvolutionaryLoop}
\end{center}
\end{figure}

\section{Mutation Operators}
\label{section:MutationOperators}

\subsection{Classical Mutation in Genetic Programming}

Classical \emph{genetic programming}\index{Genetic programming} (GP) represents expressions as trees and mutates them by subtree operations:
\begin{itemize}
\item \textbf{Subtree mutation}: replace a randomly chosen subtree with a randomly generated one.
\item \textbf{Point mutation}: replace an operator or terminal at a randomly chosen node.
\item \textbf{Crossover}: swap subtrees between two parent expressions.
\end{itemize}
The three operations only mean something once the representation is in view.
An expression is stored not as a string of characters but as a tree whose internal nodes are operators and whose leaves are variables or numeric constants: \code{x0 * exp(-x1)} is a tree with $\times$ at the root, the leaf \code{x0} hanging off one side, and the small subtree \code{exp(-x1)} hanging off the other.
A \emph{subtree} is any node together with everything beneath it, so \code{exp(-x1)}, \code{-x1}, and the bare leaf \code{x0} are all subtrees of that expression, and each operator above is a way of swapping such pieces in and out.
On that tree, a subtree mutation might replace \code{exp(-x1)} with a freshly generated \code{x1 + x2}, giving \code{x0 * (x1 + x2)}; a point mutation might change the root $\times$ to $+$, giving \code{x0 + exp(-x1)}; and a crossover with a second parent \code{x2 / (1 + x1)} might graft that parent's \code{1 + x1} in place of \code{exp(-x1)}, giving \code{x0 * (1 + x1)}.
These operations are \emph{blind} as \emph{variation}: the mutation and crossover operators propose changes without regard to whether the change improves fitness, and direction enters only later, through selection.
Classical practice is not monolithic: genetic-programming variants, symbolic-regression packages, and program synthesis differ, and some add typed or guided operators; but undirected variation is the shared baseline the next section improves on.
GP is effective when the search space is smooth enough that random perturbations frequently produce useful variation.
Smooth here means that a small edit usually causes a small change in score, so a run can improve by accumulating small wins.
Expression spaces are only partly smooth in that sense: flipping one operator can turn a curve that decays into one that diverges, and the score jumps by orders of magnitude for a change of a single node.

\subsection{LLM as Mutation Operator}

A large language model as the mutation operator replaces blind or random tree operations with \emph{knowledge-guided variation}.\index{LLM mutation operator}
Instead of choosing a random subtree to replace, the LLM reads the current best expression, the data, and domain context, and proposes a specific, motivated variant.

\begin{example}
A mutation prompt:
\begin{lstlisting}[language=markdown]
The current best model is:
    expression: "x0 * exp(-x1)"
    MSE: 0.0045, complexity: 5, score: 0.0055   # score = MSE + lambda*C, lambda = 0.0002

The dataset appears to describe exponential decay with a rate that
varies between samples (column x1). Retrieved domain note:
  "Decay rates in biological systems often follow Arrhenius kinetics:
   rate ~ exp(-Ea / (k * T)), where T is temperature."

Propose one mutation of the current expression that might better
capture this behavior. Return the expression as a Python string
using variables x0, x1, x2.
\end{lstlisting}
Three ingredients in that prompt are doing the work, and none of them is available to a blind tree edit: the current expression \emph{together with its score}, so the proposer knows what it has to beat; a short reading of what the data appears to be doing; and a retrieved domain note that names a mechanism.
The LLM might return \code{"x0 * exp(-x1 / x2)"} — a motivated variation informed by domain knowledge that a blind mutation would find only by chance.
The edit is small (one division inserted) but it is not arbitrary: the Arrhenius note says the decay rate should carry a temperature in the denominator of the exponent, and \code{x2} is the column that plausibly plays that role.
A subtree mutation could produce the same expression, but only by drawing \code{x1 / x2} at random out of the space of things it might have drawn.
\end{example}

This is the advantage that separates LLM-guided evolution from classical genetic programming, including mature libraries such as PySR, the GP-based symbolic-regression baseline on the course ladder (Table~\ref{table:BaselineLadder}), whose variation operators remain essentially blind tree edits.
The LLM mutation operator carries \emph{priors} about which variations are likely to be meaningful, and it can use \emph{domain context} to make those priors specific.

\subsection{Context Engineering for the Mutator}

The quality of the LLM mutation operator is directly determined by the quality of its context.
The mutator's context should include:
\begin{enumerate}
\item The current best expression and its score.
\item The history of what has been tried and why it did or did not work.
\item Retrieved domain knowledge relevant to the data (possibly leveraging the retrieval system of Chapter~\ref{chapter:Retrieval}).
\item The population's diversity; if all survivors are variations of the same form, the mutator should be told to explore a different region.
\end{enumerate}
The second and fourth items are the ones most often left out, and both cost real budget when they are.
Without a record of what has already failed, the mutator will cheerfully re-propose an expression that was scored and rejected three generations ago, paying again to learn the same fact.
Without a diversity summary, it has no way to know that all eight survivors are minor variations on one form; a single line in the prompt reporting that fact is often enough to push the next proposal into a different region.
This is the convergence of the two spines: \emph{evaluation becomes fitness} (Spine~1) and \emph{context engineering drives the mutator} (Spine~2).
Figure~\ref{figure:TwoSpinesConverge} shows the two threads entering the loop at the two nodes they govern: the evaluation harness \emph{becomes} the fitness that scores each candidate, and context engineering \emph{drives} the mutator that proposes the next one.

\begin{figure}[htb!]
\begin{center}
%% two-spines-converge.tex — the book's conceptual climax (Chapter 10):
%% the two spines enter the evolutionary loop. House conceptual style.
%% _concept-style.tex: shared TikZ styles for the course's CONCEPTUAL block diagrams.
%% The leading underscore marks this as the one file in figures/ that is not a
%% figure; every other figures/*.tex holds exactly one figure body.
%%
%% \input this at the top of every conceptual figure (before \begin{tikzpicture})
%% so each such diagram speaks one visual language: rounded "block" nodes, stealth
%% arrows, and natural `<hue>!15' fills. Redefining these styles on each \input is
%% idempotent and harmless. The reference for the house parameters is
%% resources/STYLE-figures.md. Figure~\ref{figure:AgentLoop} is the exemplar.
%%
%% Scope: use for conceptual/relational diagrams (boxes-and-arrows). Do NOT use for
%% product mock-ups (the rig figures have their own shared language) or for
%% data/plot figures.
\tikzset{
  % the standard block node: geometry only — apply a `fill=<hue>!15' per node.
  conceptbox/.style={draw, rounded corners=4pt, minimum width=2.4cm,
                     minimum height=0.85cm, align=center, font=\small},
  % a flow arrow between blocks.
  conceptflow/.style={->, >=stealth', thick},
  % an edge/annotation label.
  conceptlabel/.style={font=\scriptsize, align=center},
}%

\begin{tikzpicture}[
    spine/.style={->, >=stealth', thick, densely dashed},
]
% --- the two spine sources (left) ---
\node[conceptbox, fill=green!20] (harness) at (0, 2.0)  {Evaluation\\Harness};
\node[conceptbox, fill=blue!15]  (context) at (0,-2.0)  {Context\\Engineering};
\node[conceptlabel, above=1pt of harness] {Spine 1 (Chapter~\ref{chapter:Evaluation})};
\node[conceptlabel, below=1pt of context] {Spine 2 (Chapters~\ref{chapter:ContextEngineering}, \ref{chapter:Retrieval})};

% --- the evolutionary loop (right): a triangle Fitness -> Select -> Mutate ---
\node[conceptbox, fill=green!20]  (fit) at (5.6, 2.0)  {Fitness};
\node[conceptbox, fill=orange!20] (sel) at (8.0, 0.0)  {Select};
\node[conceptbox, fill=blue!15]   (mut) at (5.6,-2.0)  {Mutate};

% internal loop arrows
\draw[conceptflow] (mut) -- node[conceptlabel, left]  {candidates} (fit);
\draw[conceptflow] (fit) -- node[conceptlabel, above right, pos=0.4] {scored} (sel);
\draw[conceptflow] (sel) -- node[conceptlabel, below right, pos=0.6] {survivors} (mut);

% --- the two spines feeding the loop ---
\draw[spine] (harness) -- node[conceptlabel, above] {becomes} (fit);
\draw[spine] (context) -- node[conceptlabel, below] {drives}   (mut);
\end{tikzpicture}%

\caption{The two spines converge on the evolutionary loop.
\emph{Spine~1} (evaluation, Chapter~\ref{chapter:Evaluation}) \emph{becomes} the fitness node that scores each candidate; \emph{Spine~2} (context engineering, Chapters~\ref{chapter:ContextEngineering} and~\ref{chapter:Retrieval}) \emph{drives} the mutation node that proposes the next one.
Neither is new machinery: both are foundations threaded since the opening chapters, now meeting on one artifact.}
\label{figure:TwoSpinesConverge}
\end{center}
\end{figure}

\begin{example}
To see the two spines converge on one artifact, follow a single expression through one generation.
The survivor \code{x0 * exp(-x1)} enters the round carrying the harness score of Chapter~\ref{chapter:Evaluation}: $\mathrm{MSE} = 0.0045$, complexity $C = 5$, and composite $\Score = \mathrm{MSE} + \lambda C = 0.0055$ at $\lambda = 0.0002$.
The generate--evaluate--select loop of Chapter~\ref{chapter:Orchestration} selects it as an elite and passes it, together with that very score and a retrieved domain note on Arrhenius kinetics, into the mutation prompt shown above.
The mutator returns the motivated variant \code{x0 * exp(-x1 / x2)}.
The \emph{same} harness now scores the variant: $\mathrm{MSE} = 0.0021$, complexity $C = 7$, and composite $\Score = 0.0021 + 0.0002 \cdot 7 = 0.0035$.
Because $0.0035 < 0.0055$, selection keeps the variant, and the next mutation prompt will carry \emph{its} score in turn.
No new machinery appears in this trace: the number that drives selection is the Chapter~\ref{chapter:Evaluation} evaluation score, now used to \emph{drive} search rather than to \emph{compare} methods.
That reuse --- one harness, one score, two roles --- is what ``evaluation becomes fitness'' means concretely.
\end{example}

\section{Selection and Population Management}
\label{section:Selection}

\subsection{Elitism}

\emph{Elitism}\index{Elitism} is the practice of carrying the top-$k$ candidates from one generation to the next, regardless of what mutations are generated.
Without elitism, a mutation round can accidentally reduce the best score in the population, a random regression.
With elitism, the best score is monotonically non-increasing across generations.
In plain terms, the running best can improve or stay flat but can never get worse, because the incumbent champion is copied into the next generation untouched and can only be displaced by a candidate that actually beat it.
The guarantee is cheap and worth having for a practical reason as much as a theoretical one: it makes the leaderboard curve readable, since a flat stretch is then genuine stagnation rather than a good candidate that fell out of the population by accident.

\subsection{Selection Strategies}

Two selection strategies are commonly used.

\begin{enumerate}
\item \textbf{Tournament selection.}\index{Tournament selection}
Draw $s$ candidates at random from the population; return the best among them.
The parameter $s$ (the tournament size) controls selection pressure: a large $s$ selects the best more aggressively; a small $s$ preserves more diversity.
With a population of twenty and $s = 3$, a middling candidate wins its tournament every so often, simply because the two it was drawn against happened to be worse, and that is the point: tournament selection keeps material alive that a strict ranking would have thrown away.

\item \textbf{Top-$k$ selection.}\index{Top-$k$ selection}
Rank all candidates by fitness and retain the top $k$.
This is the simplest strategy and is the default in the course.
It is also deterministic, which makes a run easier to reproduce and to debug, at the price of the diversity that tournament randomness buys.
\end{enumerate}

\subsection{Population Collapse}

\emph{Population collapse}\index{Population collapse} occurs when the population converges to a single candidate (or a cluster of nearly identical candidates), eliminating diversity.
Diversity is the fuel of evolution: without variation, the mutation operator has nothing to vary on, and the search stagnates at a local optimum.
Collapse has a recognizable signature once you know to look for it.
The best score stops moving, and the survivor list printed each generation reads as several cosmetic rewrites of one formula: \code{x0 * exp(-x1)}, \code{1.0 * x0 * exp(-x1)}, \code{x0 / exp(x1)}, and so on.
Mutating any member of such a list lands back in the same neighborhood, so the run keeps paying for evaluations that cannot teach it anything it does not already know.

\begin{remark}
Population collapse is one of the most common failure modes in early implementations.
If the leaderboard score of the evolutionary run does not improve over the GES loop of Chapter~\ref{chapter:Orchestration}, the most likely cause is collapse: the population converged early, and the mutation operator is spinning on a single form without finding anything new.
\end{remark}

\section{Fitness Design}
\label{section:FitnessDesign}

The fitness function is the evaluation harness.
In the evolutionary setting, the harness is the target of a persistent, tireless optimizer; and every weakness in the harness is an attack surface.
The difference in kind is worth stating plainly.
A researcher running a handful of fits by hand will never notice that the harness quietly drops rows where the expression returns \code{NaN} before averaging the squared errors; an optimizer running ten thousand candidates will find that hole within a few generations and settle into it, because an expression that fails on the two hardest validation points now scores better than one that is merely accurate everywhere.
Nothing about that is malice, and the fix is not to ask the search to behave: the fix is to close the hole, because the search is doing exactly what the score told it to do.

\begin{remark}
The attack surface widens sharply when the fitness evaluation itself runs in a tool-enabled environment.
As of early 2026, capable models have been observed to recognize that they are being evaluated, enumerate likely benchmarks by name, and --- given web search and code execution --- locate a benchmark's own answer materials, turning ``solve the task'' into ``solve the evaluation'' (Anthropic, 2026).
The defenses are the ones the leakage discussion of Chapter~\ref{chapter:Evaluation} already prescribes: run the fitness function in a sandbox with no route to the sealed held-out data or the open web, and treat honest-failure reporting as part of the harness contract rather than an afterthought.
\end{remark}

\subsection{Multi-Objective Fitness}

The course harness uses the composite score $\Score = \mathrm{MSE} + \lambda \cdot C$ defined in~\eqref{equation:CompositeScore}.
The composite folds two things a researcher genuinely cares about into the one number a search can minimize: $\mathrm{MSE}$ measures how well the expression reproduces held-out data, the complexity $C$ counts the nodes in its tree, and $\lambda$ sets the exchange rate between them.
At $\lambda = 0.0002$, adding one node to an expression must buy an MSE improvement of at least $0.0002$ to be worth it, so $\lambda$ is best read as the price of a node rather than as an abstract regularization weight.
In the evolutionary setting, a well-designed $\lambda$ prevents the optimizer from discovering that a trivially complex, trivially fitted expression scores as well as a genuinely discovered one.

For long evolutionary runs, the optimizer often finds unexpected ways to minimize this score.
Common exploits:
\begin{enumerate}
\item An expression of zero complexity: a constant fitted to the training data.
If $\mathrm{MSE}$ is computed on the validation split, a global constant performs worse than a simple model on diverse points; but if $\lambda$ is too large relative to the MSE gap, the complexity reduction can dominate.
\item Expressions that overfit the validation split via extreme complexity: $10^6$-order polynomials.
The parsimony term must scale to prevent this.
\end{enumerate}

\subsection{Fitness Landscapes and Local Optima}

The \emph{fitness landscape}\index{Fitness landscape} is the surface of fitness values over the space of artifacts.
The metaphor is worth taking literally for a moment: picture every candidate expression as a point on a map, with two expressions adjacent when one mutation separates them, and the elevation at each point given by its score.
Because this course minimizes, the search is walking \emph{downhill}, and a \emph{basin} is a valley of that landscape: a region from which every downhill step leads to the same low point.
A local optimum is a point where all nearby mutations produce worse fitness, even though a distant point might be better.
The evolutionary algorithm can become trapped at local optima if selection pressure is too high (population collapse) or mutation step size is too small (unable to escape the basin).
A concrete trap in symbolic regression: a run that has locked onto a polynomial form keeps improving its score for many generations by adding one term at a time, while the process that actually generated the data is an exponential that no sequence of small polynomial edits ever reaches.
The run looks healthy on the leaderboard, and it is nonetheless in the wrong valley.

Strategies for avoiding local optima:
\begin{itemize}
\item Run multiple independent populations (islands) and periodically share the best candidates between them.
\item Allow the mutator to make larger jumps; propose entirely new functional forms rather than just small edits to existing ones.
\item Maintain explicit diversity pressure: penalize candidates that are too similar to existing survivors.
\end{itemize}

\section{Diversity and Exploration}
\label{section:Diversity}

\subsection{Islands}

An \emph{island model}\index{Island model} runs several independent populations in parallel, each exploring a different region of the search space.
Periodically, the best candidates from each island \emph{migrate} to other islands, injecting diversity.
The island model prevents all populations from converging to the same local optimum.

\begin{example}
A four-island setup in pseudocode:
\begin{lstlisting}[language=python]
# Initialize four independent populations
islands = [Population(seed=i) for i in range(4)]

for generation in range(100):
    # Each island evolves independently
    for island in islands:
        island.evolve(fitness_fn)

    # Every 10 generations, share the best from each island
    if generation % 10 == 0:
        best_per_island = [island.best() for island in islands]
        for island in islands:
            island.inject(best_per_island)
\end{lstlisting}
Two details in that sketch carry the design.
Each island is constructed with a different seed, so the four populations start in different regions of the landscape instead of being four copies of the same run.
And migration happens every ten generations, not every generation: migrate too often and the islands merge back into one population, which is the situation the model exists to prevent, while migrating too rarely spends parallel compute on four populations that have each stalled in their own valley.
\end{example}

\subsection{Explore vs.\ Exploit}

The \emph{explore–exploit tradeoff}\index{Explore--exploit tradeoff} is the fundamental tension in any search:
\begin{itemize}
\item \textbf{Exploit}: mutate the best survivors to find small improvements.
\item \textbf{Explore}: generate new candidates far from the current population to discover new regions.
\end{itemize}
Too much exploitation leads to premature convergence; too much exploration wastes evaluation budget on low-quality candidates.
The tension is familiar well outside search: an experimentalist who only ever refines the assay that already works will never find the better assay, while one who only ever tries new assays never gets a clean measurement out of any of them.

A practical heuristic: maintain a diverse population (island model or explicit diversity penalty), use a proposer that mixes small mutations (exploit) with occasional large jumps (explore), and increase the exploration rate if the best score has not improved for several consecutive generations.

\section{The Lineage: FunSearch, AlphaEvolve, OpenEvolve}
\label{section:Lineage}

The agentic evolutionary framework taught in this course has a direct intellectual lineage.

\subsection{FunSearch}

\emph{FunSearch}\index{FunSearch} (Romera-Paredes et al., DeepMind; \emph{Nature}, 2024, online 2023) is an LLM-guided evolutionary search over \emph{programs} (Python functions), scored by an automatic evaluator.
The system evolved improved constructions for combinatorics problems (finding larger cap sets in $\mathbb{F}_3^8$) and better heuristics for online bin packing.
A \emph{cap set} in that space is a collection of points containing no three that lie on a line, and the long-standing question is how large such a collection can be; the point for a reader outside combinatorics is only that the answer is a discrete object a short program can construct and a short script can check, which is exactly the shape a fitness function needs.
Bin packing is the everyday counterpart: items arrive one at a time and must be placed into bins without knowing what comes next, and the evolved artifact is the rule deciding where each item goes.

The template introduced by FunSearch is:
\begin{itemize}
\item \textbf{Artifact}: a Python function.
\item \textbf{Mutation}: the LLM rewrites the function, guided by examples of high-scoring past versions.
\item \textbf{Fitness}: an automatic correctness and quality evaluator specific to the problem.
\item \textbf{Infrastructure}: a population of functions, maintained and sampled across many parallel evaluations.
\end{itemize}
This is precisely the structure implemented in the capstone.

\subsection{AlphaEvolve}

\emph{AlphaEvolve}\index{AlphaEvolve} (DeepMind, 2025) scales the FunSearch template to evolve full algorithms using Gemini models.
Results include improved matrix-multiplication procedures (beating the Strassen algorithm in certain settings) and real scheduling gains in data-center resource allocation.
Strassen's algorithm is the classical trick for multiplying two matrices using fewer multiplications than the textbook definition requires, and shaving further multiplications off it has been a target for decades; a system that improves on it, even in restricted settings, is competing on a problem where human effort has been concentrated for a long time.
AlphaEvolve demonstrates that the approach scales: with sufficiently capable models and evaluators, the method finds results competitive with specialized mathematical research.

\subsection{OpenEvolve}

\emph{OpenEvolve}\index{OpenEvolve} is an open-source reimplementation of the AlphaEvolve approach.
It is the primary reference implementation for the capstone: a hackable codebase that exposes the mutation operator, fitness function, and population management as modifiable components.
The studio session at the pivot walks through OpenEvolve on a Tier-1 symbolic regression dataset, tracing each component to its corresponding concept in these notes.

\subsection{The AI Scientist (Optional)}

\emph{The AI Scientist} (Sakana AI, 2024) extends the evolutionary template to the full research cycle: automated literature review, hypothesis generation, experiment execution, and paper writing.
It is more agentic-research than evolutionary-search in structure, and is included as a glimpse of the long-run trajectory rather than as a primary reference.

\section{Symbolic Regression as the Teaching Vehicle}
\label{section:SymbolicRegression}

\subsection{The Model Discovery Problem}

\emph{Symbolic regression}\index{Symbolic regression} (SR) is the task of discovering a mathematical expression $\hat{m}$ that describes the relationship between a set of features $X$ and a target $y$.
Unlike conventional regression, the functional form of $\hat{m}$ is not fixed in advance; the search must explore a large, discrete space of possible expressions.
The contrast is worth making concrete.
Fitting $y = \beta_0 + \beta_1 x_1 + \beta_2 x_2$ estimates three numbers with the shape of the model handed to you in advance; symbolic regression asks for the shape as well, so the answer can come back as $y = x_0 e^{-x_1 / x_2}$ without anyone having thought to propose a ratio inside an exponential.
That freedom is what makes the space discrete and enormous: the number of syntactically valid expressions grows faster than exponentially in the node budget, so exhaustive enumeration is off the table even for expressions a human would call short, and every method in this chapter is a different way of not enumerating.

\begin{definition}
\label{definition:SymbolicRegression}
Given a dataset $\mathcal{D} = \{(x_i, y_i)\}_{i=1}^{N}$ and a space of expressions $\Artifacts$, \emph{symbolic regression} seeks
\begin{equation}
\hat{m}^* = \operatorname*{arg\;min}_{a \in \Artifacts} \Score(a, \mathcal{D}_{\mathrm{val}})
\end{equation}
where $\Score$ is the composite score of~\eqref{equation:CompositeScore} and $\mathcal{D}_{\mathrm{val}}$ is the validation (fitness) split of the harness (Definition~\ref{definition:EvalHarness}); the discovered $\hat{m}^*$ is then reported once against the untouched final-test split.
\end{definition}

In words: search the space of expressions, return the one that scores best on the validation split, then report that single winner once against a test split no candidate was ever scored on.
The two splits do different jobs, and blurring them is the most common way an impressive-looking discovery evaporates.
The validation split is consumed by the search, so its score is optimistic by construction: thousands of candidates were filtered through it, and the winner is partly a winner because of the particular noise in those rows.
The untouched test split is the only honest estimate of how the discovered model behaves on data that had no hand in shaping it, which is why it is looked at once, at the end, and never used to choose between candidates.
This is the same discipline as pre-registering an analysis, transposed to a search that runs thousands of times.

A note on convention: the course \emph{minimizes} throughout, so the composite score is a \emph{loss} (lower is better); we nonetheless call the same quantity \emph{fitness} where it drives selection, so ``worse fitness'' means a higher score.
Where evolutionary computation would have fitness increase with quality, read the usage here as selecting for lower loss.

SR is the teaching vehicle for the course because it is:
\begin{enumerate}
\item \textbf{Automatically evaluable.} The fitness is a number computed without human judgment.
\item \textbf{Domain-portable.} A biologist can discover enzyme kinetics; a physicist can discover force laws; an economist can discover demand curves.
Each of those is the same run with a different data file: the biologist's columns are substrate concentration and reaction velocity, and a winning expression may be the Michaelis--Menten form $v = V_{\max} S / (K_m + S)$; the physicist's columns are separation and measured force; the economist's are price, income, and quantity demanded.
The domain changes; the framework does not.
\item \textbf{Interpretable.} The output is a symbolic expression, not a neural network.
The researcher can read and critique the discovered model.
\item \textbf{Has established baselines.} Linear regression and PySR (a mature GP-based SR library) provide calibration points against which the evolutionary agent must compete.
\end{enumerate}

\subsection{The Baseline Ladder}

The course maintains a \emph{baseline ladder}\index{Baseline ladder}: a sequence of methods of increasing sophistication against which each researcher's system is compared.
Beating the each rung is not required; being comparable to it, on a well-designed harness, is.

\begin{table}[htb!]
\begin{center}
\begin{tabular}{lll}
\toprule
\textbf{Rung} & \textbf{Method} & \textbf{Challenge} \\
\midrule
1 & Linear regression & Challenge~1 \\
2 & Research pipeline (LLM proposals) & Challenge~2 \\
2b & PySR (classical GP-based SR) & Challenge~2 benchmark \\
3 & Retrieval-grounded pipeline & Challenge~3 \\
4 & Orchestrated multi-proposer & Challenge~4 \\
5 & Evolutionary framework & Capstone \\
\bottomrule
\end{tabular}
\caption{The baseline ladder for the course.
Each rung represents a more capable method; the leaderboard shows all entries simultaneously.
A stronger method is expected to score lower (better) than weaker methods on the same dataset.}
\label{table:BaselineLadder}
\end{center}
\end{table}

\section{Cost-Aware Search}
\label{section:CostAwareSearch}

Every fitness evaluation and mutation is an API call that costs money and takes time.
A long evolutionary run on a high-cost model can be expensive; a run on a cheap model may produce inferior mutations that the run cannot recover from.

\begin{definition}
The \emph{cost-adjusted fitness} of an artifact $a$ is
\begin{equation}
\tilde{\Fitness}(a) = \Fitness(a) + \alpha \cdot \mathrm{cost}(a) ,
\label{equation:CostAdjustedFitness}
\end{equation}
where $\mathrm{cost}(a)$ is the \emph{marginal} cost (in dollars of API usage) of producing and evaluating $a$, and $\alpha > 0$ converts that cost to the scale of the fitness score.
\end{definition}

In words, an artifact is charged for what it cost to obtain, in the same units as its score, so a candidate is kept only when its quality justifies its price.
Suppose $\alpha$ is set so that one dollar of API usage is worth $0.001$ of score.
A variant that improves the composite score by $0.0005$ but took ten dollars of proposer calls to find is then charged $0.01$ and ranks as a loss rather than a gain, while the same improvement obtained for a few cents survives comfortably.

Using the \emph{marginal} cost $\mathrm{cost}(a)$, rather than a run's total cost, is what makes the adjustment discriminating: a cost shared equally by every artifact would shift all fitnesses by the same constant and could not reorder candidates within a run.
With $\mathrm{cost}(a)$, two candidates of equal $\Fitness$ are separated in favor of the one that was cheaper to obtain — the mutation drawn from a cheaper proposer model, or the expression that reached its score in fewer evaluation calls.
This connects the capstone to the multi-model benchmarking of Chapter~\ref{chapter:Orchestration}: a cheaper proposer that achieves comparable mutations is preferred by cost-adjusted fitness, even when its per-mutation quality is slightly lower.

\begin{remark}
Like the parsimony weight $\lambda$ (Section~\ref{subsection:MultiObjectiveScoring}), $\alpha$ is a calibrated design choice, not a universal constant.
Its units are fitness per dollar, so its scale depends on both the MSE scale of the dataset and the dollar cost of a typical mutation.
A workable default fixes $\alpha$ so that halving an artifact's cost is worth the same as a small, stated fraction of the MSE gap between adjacent baselines.
For research purposes the primary leaderboard uses unadjusted fitness so that discovery quality is the dominant signal; the cost analysis is reported separately, enabling the reader to understand the resource requirements of each method.
More fundamentally, generation cost is a property of the \emph{path} that produced an artifact, not of the artifact itself, the same expression can be reached cheaply or expensively.
Folding cost into an artifact's fitness can distort survivor selection.
Read the cost-adjusted scalar as scoring a search \emph{policy} rather than an artifact, and prefer a run-level budget or an explicit Pareto view (quality, complexity, dollars, wall-clock) when quality and cost cannot be honestly collapsed into one number.
\end{remark}

\section{What You Have Built}
\label{section:WhatYouHaveBuilt}

The chapters of these notes form one argument.
A researcher who can operate a rig, build a loop, engineer context, evaluate honestly, and orchestrate reliably has already assembled every component of an evolutionary search.
The mutation operator of this chapter is the one part that was missing, and it reuses the proposer, the harness, and the orchestration that came before.
The capstone is therefore not a new topic; it is these components, finally pointed at your own research problem, and the artifact they evolve is a model of your own data.

Pointing it at your own problem is, concretely, a matter of answering the five questions the tuple of Definition~\ref{definition:EvolutionaryLoop} asks.
What is the artifact you want the search to produce: an expression, a model specification, an experimental protocol, a controller, a design?
What scores one artifact automatically, with no one in the loop, and where does the data that scores it come from?
What counts as a small, meaning-preserving edit to an artifact, and what context does a proposer need in order to make a good one?
Which survivors get to be parents, and how many?
And where does the first population come from?
A problem for which the second question has no honest answer is not yet ready for this machinery, and reaching that verdict early is itself worth the exercise: it says the next piece of work is building the measurement, not the search.

\section*{Further Reading}

\begin{small}
\begin{enumerate}
\item Romera-Paredes, B., et al., ``Mathematical Discoveries from Program Search with Large Language Models,'' \emph{Nature}, 2024 (online December 2023; DOI:10.1038/s41586-023-06924-6) — the FunSearch paper; the primary template for the capstone.
\item DeepMind, ``AlphaEvolve: A Gemini-Powered Coding Agent for Designing Advanced Algorithms,'' technical report, 2025 — the scaled-up successor to FunSearch.
\item Lu, C., et al., ``The AI Scientist: Towards Fully Automated Open-Ended Scientific Discovery,'' Sakana AI, 2024 (arXiv:2408.06292) — a related but distinct end-to-end research-automation workflow, referenced in this chapter's lineage discussion.
\item Shinn, N., et al., ``Reflexion: Language Agents with Verbal Reinforcement Learning,'' \emph{NeurIPS}, 2023 (arXiv:2303.11366) — an evaluation signal turned into \emph{linguistic} feedback for the next attempt. Reflexion accumulates that feedback across repeated attempts at one task; the mutation operator here compares across a population of independent candidates --- related but distinct uses of language as a feedback signal.
\item Yao, S., et al., ``Tree of Thoughts: Deliberate Problem Solving with Large Language Models,'' \emph{NeurIPS}, 2023 (arXiv:2305.10601) — search over intermediate reasoning steps \emph{within} a single model's deliberation; the internal-search foil to this chapter's external, population-based search.
\item OpenEvolve repository — open-source AlphaEvolve implementation; see \code{benchmarks/README.md} for the course-specific setup.
\item Koza, J. R., \emph{Genetic Programming}, MIT Press, 1992 — the classical reference for GP and symbolic regression; Chapter~6 covers the baseline methods against which the course agents compete.
\item Schmidt, M., and Lipson, H., ``Distilling Free-Form Natural Laws from Experimental Data,'' \emph{Science}, 2009 (DOI:10.1126/science.1165893) — the origin of modern SR; PySR's scientific lineage traces here.
\item \texttt{projects/capstone/evolution-pivot.md} in the course repository — the pivot lecture notes, including the OpenEvolve studio walkthrough.
\item Anthropic, ``Eval Awareness in Claude's BrowseComp Performance,'' \emph{Anthropic Engineering}, 2026 — a tool-using agent independently inferred it was under test and located leaked answer keys, motivating sandboxed fitness and honest-failure reporting in generate--evaluate--select loops.
\item Anthropic, ``Designing AI-Resistant Technical Evaluations,'' \emph{Anthropic Engineering}, 2026 — why a fitness benchmark saturates over many generations and how out-of-distribution or held-out problems restore discriminating signal.
\end{enumerate}
\end{small}
% EVOLVE-BLOCK-END
